%
% quarry-lux-search-applet-prd.tex
%
% Product requirements for a quarry lux applet: search across every
% registered collection, and a read-only view of daemon status, reachable
% from the lux menu bar whenever quarryd is running -- with no dependency on
% any particular Claude Code session being open. Written from the
% product-manager altitude -- what the feature must do and for whom, not
% how it is implemented. Use cases and functional requirements are traced
% to each other so a reviewer can check coverage in both directions.
% Architecture (the QuarrySearchSubsystem composition, the lux board/scene
% wire shapes) is deliberately out of this document; it follows, once this
% is ratified, per docs/WORKFLOW.md's design-before-implementation gate.
%
% Companion document: quarry-lux-control-panel-prd.tex, a separate,
% session-bound applet controlling this repo's own capture configuration.
% The two are deliberately different applet species (DES-063, lux repo)
% and are scoped as two independent epics, not one.
%

\documentclass[a4paper,11pt]{report}
\usepackage[margin=1in]{geometry}
\usepackage[colorlinks=true,linkcolor=blue,citecolor=blue,urlcolor=blue]{hyperref}
\usepackage{enumitem}
\usepackage{longtable}
\usepackage{booktabs}
\usepackage{parskip}

\newcommand{\uc}[1]{\textbf{UC-#1}}
\newcommand{\fr}[1]{\textbf{FR-#1}}
\newcommand{\nfr}[1]{\textbf{NFR-#1}}
\newcommand{\pone}{\textsc{[p1]}}
\newcommand{\ptwo}{\textsc{[p2]}}

\begin{document}

\title{Quarry Lux Search Applet: Product Requirements}
\author{Punt Labs}
\date{August 2026}
\hypersetup{
  pdftitle={Quarry Lux Search Applet: Product Requirements},
  pdfauthor={Punt Labs},
  pdfsubject={Product Requirements Document}
}
\maketitle
\tableofcontents

\chapter{Product Requirements}

\section{Overview}
\label{sec:overview}

Quarry today has one graphical entry point: \texttt{quarry-menubar}, a
native macOS \texttt{MenuBarExtra} app that connects to \texttt{quarryd}
over its resolved connection (local TLS or a configured remote server) and
lets a user search and read indexed content. This document specifies a
second entry point: a lux applet, reachable from the shared lux menu bar
whenever \texttt{quarryd} and \texttt{luxd} are both running, offering the
same core search-and-read experience to any lux user regardless of
platform or whether \texttt{quarry-menubar} is installed.

This is not a replacement for \texttt{quarry-menubar} and does not reduce
its scope. It is an additional surface for the same backend, matching how
\texttt{voxd}'s music player already gives vox a lux-native control surface
alongside vox's own CLI and MCP tools.

The applet is machine-global, not session-scoped: it is composed directly
into \texttt{quarryd}'s own process at daemon startup (the same shape as
\texttt{voxd}'s \texttt{MusicPlayerSubsystem}), not spawned per Claude Code
session. This is a product decision, not an implementation detail --- a
search across every registered collection, and the health of the daemon
serving them, are both facts about the machine, not about any one editor
session, and the applet's availability should track the daemon's uptime,
not a session's.

\subsection{Priorities and phasing}

\pone{} is a working search-and-read applet, at feature parity with
\texttt{quarry-menubar}'s core loop (query, browse results, read full
content) reachable from the lux menu bar, plus a read-only status view.
\ptwo{} is polish and parity items that improve the experience without
being required to call the feature done: richer status detail, performance
tuning for large result sets, and any lux-side affordance
\texttt{quarry-menubar} has that this document's \pone{} scope omits.

\subsection{Goals}

\begin{itemize}
  \item Let any lux user search every collection quarry has indexed,
    from the lux menu bar, without installing or running
    \texttt{quarry-menubar}.
  \item Make reading a matched document's full content --- not just the
    search snippet --- the primary experience, matching what
    \texttt{quarry-menubar} already proved is the right shape (\uc{2},
    \uc{3}).
  \item Surface enough daemon health at a glance that a user can tell
    whether a failed or empty search is quarry's fault or their query's
    fault, without leaving lux to run \texttt{quarry doctor}.
  \item Reuse the \emph{contract} quarry's existing HTTP API already
    defines (\texttt{/v1/search}, \texttt{/v1/show},
    \texttt{/v1/collections}, \texttt{/v1/status}) --- the same request
    shapes, the same response fields --- as-is where it already serves the
    need; extending the contract is in scope only where a requirement
    below cannot be met without it. \textbf{Revised (Chapter 2,
    \S2.4):} this applet reaches that contract's underlying service
    layer \emph{in-process}, not over HTTP loopback, since it runs
    inside \texttt{quarryd} itself rather than as an external client ---
    the goal is contract parity with the HTTP API, not literally issuing
    HTTP requests to reach it.
\end{itemize}

\subsection{Non-Goals}

\begin{itemize}
  \item Writing to quarry: no ingest, no \texttt{remember}, no delete, no
    collection management. This applet reads; it does not mutate quarry's
    index. (Distinguish from the companion control-panel applet, which
    does mutate --- but only this repo's own capture configuration, never
    the index.)
  \item Any daemon lifecycle control (restart, service management) or
    connection/remote-config management (API keys, TLS trust). That stays
    \texttt{quarry-menubar}'s job; see \S~\ref{sec:noncontrol} below.
  \item Replicating \texttt{quarry-menubar}'s keyboard-first, hotkey-summoned
    UX exactly. Lux's own menu-bar interaction model (a click raises a
    frame) is the entry point here; matching \texttt{quarry-menubar}'s
    specific hotkey behavior is out of scope.
  \item Any table/filter/type-combo chrome in the style of the beads
    board (\texttt{apps/beads\_board.py}'s live-chrome \texttt{RenderTableRequest}).
    Considered and explicitly deferred to a later iteration once the
    content-first \pone{} scope ships and is validated.
\end{itemize}

\label{sec:noncontrol}
\textbf{Why daemon control stays out.} A prior draft of this feature
combined search, status, \emph{and} daemon control (restart, database
switching, sync-now, remote-serving toggles) into one applet. That was
corrected: \texttt{quarry-menubar} already owns daemon-level status and
control in its native surface, and duplicating that control surface in lux
would create two places a user could issue conflicting daemon commands
from, with no obvious source of truth. This applet's status section is
read-only by design.

\subsection{Personas}

\begin{description}
  \item[The lux-first user] Runs \texttt{luxd} as their primary always-on
    surface across several Punt Labs tools (vox's music player, biff's
    presence) and wants quarry search available the same way, without a
    second app to install or a dock icon to manage. Primary persona for
    \pone{}.
  \item[The cross-platform user] Uses lux on a platform
    \texttt{quarry-menubar} does not target (\texttt{quarry-menubar} is
    macOS-only \texttt{MenuBarExtra}; lux's Display renderer is not
    platform-restricted the same way). This applet is their only
    graphical way to search quarry.
  \item[The quarry-menubar user, occasionally] Has \texttt{quarry-menubar}
    installed and generally prefers it, but is at a machine or in a
    context where lux is already open and quarry-menubar is not running;
    reaches for the lux applet as a fallback rather than launching a
    second app for one query.
\end{description}

\textbf{Evidence status.} No usage data or interviews back the
cross-platform and occasional-user personas above --- \texttt{quarry-menubar}
already exists and plausibly serves the primary user base today, so this
document ships on an assertion, not a validated need. This is treated
explicitly as a validation bet, not a settled fact: \S~\ref{sec:metric}
defines what would confirm or falsify it, and \pone{} is scoped to be
cheap enough to ship and observe before committing to \ptwo{}.

\section{Use Cases}

Each use case names the actor, the trigger, the expected flow, and what
``done well'' looks like. Functional requirements in \S\ref{sec:fr} cite
the use cases they satisfy.

\begin{description}

\item[\uc{1} Open the applet from the lux menu. \pone{}] \emph{Actor: any
  lux user.} \texttt{quarryd} is running; the user clicks the ``Quarry''
  entry in the lux menu bar. \emph{Done well:} the applet's frame raises
  immediately (DES-063's raise-first principle) --- the user sees a search
  field within roughly the same latency lux-beads already demonstrates for
  its own click-to-visible-response, not a spinner while the daemon is
  contacted for the first time.

\item[\uc{2} Search and scan results. \pone{}] \emph{Actor: any lux user.}
  The user types a query into the search field. \emph{Done well:} results
  appear grouped in a way that lets the user scan quickly (matching
  \texttt{quarry-menubar}'s grouping by source format), each row compact
  enough to be a picker, not a destination --- the user should feel like
  they are choosing which document to read, not already reading it.

\item[\uc{3} Read a result's full content. \pone{}] \emph{Actor: any lux
  user.} The user selects a result row. \emph{Done well:} the applet shows
  the \emph{full page} the result came from (via the same backing call as
  \texttt{quarry-menubar}'s \texttt{/v1/show}), not merely the search
  snippet already visible in the list --- reading the content is the
  point of the feature, and the list exists only to get the user here
  quickly. Code content is legible (a monospaced or otherwise
  code-appropriate rendering), matching the intent of
  \texttt{quarry-menubar}'s syntax highlighting even if lux's rendering
  primitives cannot reproduce it exactly.

\item[\uc{4} Return to the list from a result. \pone{}] \emph{Actor: any
  lux user.} The user is reading a result's content and wants to see the
  results list again --- to pick a different result, or refine the
  query. \emph{Done well:} one action returns to the list without losing
  the search query or re-running the search.

\item[\uc{5} Copy content out of the applet. \pone{}] \emph{Actor: any lux
  user.} The user is reading a result and wants to paste its content
  elsewhere. \emph{Done well:} the full content just read is copyable in
  one action, matching \texttt{quarry-menubar}'s ``Copy'' affordance.

\item[\uc{6} Search returns nothing. \pone{}] \emph{Actor: any lux user.}
  The query matches no indexed content. \emph{Done well:} the applet says
  so plainly (no results for this query) rather than showing an empty list
  indistinguishable from ``still loading'' or a silent failure.

\item[\uc{7} Search while the daemon is unreachable. \pone{}] \emph{Actor:
  any lux user.} \texttt{quarryd} is down, unreachable, or the applet's
  own connection to it has faulted. \emph{Done well:} the applet says the
  daemon is unreachable, distinctly from ``no results,'' so the user knows
  to check quarry's own health rather than assume their query was bad.
  (See also \uc{9}.)

\item[\uc{8} Reconnect after a lux restart. \pone{}] \emph{Actor: any lux
  user.} \texttt{luxd} restarts (or the applet's own connection drops and
  recovers) while \texttt{quarryd} keeps running. \emph{Done well:} the
  applet's menu entry and any held content reappear automatically once the
  connection is back, matching the register-fresh behavior
  \texttt{voxd}'s music player already demonstrates on Hub reconnect ---
  the user does not have to notice the outage or take any recovery action.

\item[\uc{9} Check daemon status without searching. \pone{}] \emph{Actor:
  any lux user.} The user opens the applet specifically to check whether
  quarry is healthy --- daemon up, index healthy, collections in sync ---
  without intending to search anything. \emph{Done well:} a status view is
  reachable from the same applet without a separate program or menu
  entry, and it is read-only: nothing in this view lets the user change
  daemon state (see \S~\ref{sec:noncontrol}).

\item[\uc{10} Filter to one collection. \ptwo{}] \emph{Actor: any lux
  user.} The user wants results from only one indexed collection, not
  every collection quarry knows about --- matching
  \texttt{quarry-menubar}'s collection picker. \emph{Done well:} the
  filter is visible and changeable from the search view itself, not a
  separate settings screen, and persists for the remainder of the applet
  session (until lux or quarryd restarts).

\item[\uc{11} Scroll or page through many results. \ptwo{}] \emph{Actor:
  any lux user.} A query matches more results than fit on screen at once.
  \emph{Done well:} the user can reach every result, not just the first
  page's worth, without the applet becoming sluggish as the result count
  grows.

\end{description}

\section{Functional Requirements}
\label{sec:fr}

Requirements are grouped by capability area. Each cites the use case(s) it
exists to satisfy and carries the same \pone{}/\ptwo{} phase tag.

\subsection{Applet lifecycle}

\begin{description}
  \item[\fr{1} \pone{}] The applet shall be available as a single menu
    entry in the lux menu bar whenever \texttt{quarryd} is running, with
    no dependency on any Claude Code session being open. (\uc{1})
  \item[\fr{2} \pone{}] Clicking the applet's menu entry shall raise its
    frame if it is already open, and open it with a visible search field
    if it is not, before any query has been issued. (\uc{1})
  \item[\fr{3} \pone{}] The applet shall re-register its menu entry and
    restore whatever it was last showing automatically on every fresh
    connection to \texttt{luxd} (including after a \texttt{luxd} restart),
    with no user action required. (\uc{8})
\end{description}

\subsection{Search}

\begin{description}
  \item[\fr{4} \pone{}] The applet shall let the user enter a free-text
    query and issue a search against quarry's indexed content. (\uc{2})
  \item[\fr{5} \pone{}] Search results shall be grouped in a way that
    supports quick scanning (at minimum, by source format, matching
    \texttt{quarry-menubar}'s existing grouping), with each row compact
    enough to show document identity and a short excerpt, not full
    content. (\uc{2})
  \item[\fr{6} \pone{}] A query that matches nothing shall produce an
    explicit ``no results'' state, distinguishable from a loading state
    and from a daemon-unreachable state. (\uc{6})
  \item[\fr{7} \ptwo{}] The user shall be able to restrict search to one
    collection, selectable from the search view. (\uc{10})
  \item[\fr{8} \ptwo{}] The applet shall remain responsive as the result
    count grows past what fits on one screen (pagination, virtualized
    scrolling, or an equivalent), without a hard cap that silently drops
    results below the fold. (\uc{11})
\end{description}

\subsection{Reading content}

\begin{description}
  \item[\fr{9} \pone{}] Selecting a result shall fetch and display the
    \emph{full page} that result came from, not only the excerpt already
    visible in the results list. (\uc{3})
  \item[\fr{10} \pone{}] The content view shall render code-format
    results legibly (at minimum, a monospaced presentation that preserves
    whitespace and line breaks). (\uc{3})
  \item[\fr{11} \pone{}] One action shall return from the content view to
    the results list without re-issuing the search or losing the query
    text. (\uc{4})
  \item[\fr{12} \pone{}] One action shall copy the full content currently
    displayed. (\uc{5})
\end{description}

\subsection{Daemon status (read-only)}

\begin{description}
  \item[\fr{13} \pone{}] The applet shall offer a status view, reachable
    from the same applet, showing at minimum: daemon reachability, and
    whether the last search attempt succeeded or failed for a
    daemon-side reason. (\uc{9})
  \item[\fr{14} \pone{}] When the daemon is unreachable, the applet shall
    say so explicitly and distinctly from a genuine zero-result search,
    both in the search view and the status view. (\uc{7})
  \item[\fr{15} \pone{}] The status view shall contain no control that
    changes daemon state (no restart, no database switch, no
    configuration write). (\uc{9}; see \S~\ref{sec:noncontrol})
  \item[\fr{16} \ptwo{}] The status view shall list the daemon's known
    collections by name and document/chunk counts, the data
    \texttt{/v1/status} already exposes today
    (\texttt{QuarryModels.swift}'s \texttt{StatusResponse}: aggregate
    \texttt{document\_count}, \texttt{collection\_count},
    \texttt{chunk\_count} --- no per-collection breakdown). A
    per-collection last-synced indication, if wanted, is a genuinely new
    field this requirement does not assume exists; scoping it is an open
    question (\S~\ref{sec:openq-lastsync}), not something
    \texttt{quarry-menubar} already provides. (\uc{9})
\end{description}

\section{Non-Functional Requirements}

\begin{description}
  \item[\nfr{1} \pone{}] Menu-click to first visible content (a search
    field, or a held prior result) shall not be materially worse than
    lux-beads' own demonstrated click-to-visible-response latency
    (DES-063 measured lux-beads at roughly 55\,ms median for a raise, with
    the slow \texttt{bd} read moving behind that raise, not blocking it).
    The equivalent principle applies here: raise first, fill behind it.
  \item[\nfr{2} \pone{}] A daemon fault (crash, restart, or a fatal fault
    in either lux leg) shall not require the applet's host process
    (\texttt{quarryd}) to be restarted by a human --- the subsystem shall
    self-heal, matching \texttt{MusicPlayerSubsystem}'s guarded
    restart-with-backoff loop.
  \item[\nfr{3} \ptwo{}] The applet shall not measurably degrade
    \texttt{quarryd}'s own search latency for other clients (the HTTP API,
    \texttt{quarry-menubar}, the MCP server) merely by being connected and
    idle.
\end{description}

\section{Success Metric}
\label{sec:metric}

This applet exists on an unvalidated persona bet (\S~\ref{sec:overview}
personas). The falsifiable claim: within the first month after \pone{}
ships, at least some non-trivial share of searches happen through this
applet from users who also have \texttt{quarry-menubar} available ---
evidence that the applet is reaching a real, distinct moment of use
(lux already open, quarry-menubar not running) rather than sitting
unused beside an app that already does the job. If usage never
materializes among users who could have used \texttt{quarry-menubar}
instead, \ptwo{} (collection filtering, pagination, richer status) is
not justified and this applet's further investment should be
reconsidered before, not after, building it.

\section{Open Questions}

\begin{enumerate}
  \item \textbf{Resolved in Chapter 2.} Whether lux's board/scene
    primitive supports replacing a board's content in place is settled:
    yes, by pushing a fresh \texttt{RenderRequest} to the same
    \texttt{frame\_id} the results table already used --- no new board
    plumbing needed. See Chapter 2, \S\ 2.2 for the three-part proof.
  \item What quarry HTTP API surface, if any, is missing for this applet
    versus what \texttt{quarry-menubar} already calls
    (\texttt{/v1/search}, \texttt{/v1/documents}, \texttt{/v1/collections},
    \texttt{/v1/show}, \texttt{/v1/status})?
  \item Does \texttt{quarryd} taking a direct dependency on
    \texttt{punt\_lux} (as \texttt{voxd} does) have any packaging
    consequence for quarry, which ships as a CLI+MCP hybrid plugin today?
  \item \textbf{Still open; carried forward as Chapter 2, \S2.5 Decision
    4.} What, precisely, should \fr{13}'s status view show beyond
    reachability --- is FD headroom and FTS index health (already in
    \texttt{quarry doctor}'s output) the right read-only subset, or is
    that too much detail for a glance-and-close use case?
  \item \label{sec:openq-lastsync} \fr{16} originally assumed
    \texttt{quarry-menubar}'s status surface already exposes a
    per-collection last-synced timestamp; peer review found it does not
    --- \texttt{/v1/status} returns only aggregate counts
    (\texttt{StatusResponse} in \texttt{QuarryModels.swift}). Is a
    per-collection last-synced field worth adding to quarry's HTTP API
    for this applet, and if so is \texttt{quarry doctor}'s existing sync
    machinery (\texttt{compute\_sync\_plan}) the right source for it, or
    is aggregate-only status sufficient for \pone{}?
\end{enumerate}

\chapter{Technical Design Solution}

This chapter resolves the open questions raised in Chapter 1 against the
actual code of the three repositories this feature spans (lux, vox,
quarry). Nothing here is speculative where the cited source settles the
question; where it does not, the gap is named as a decision required
before implementation, not silently assumed.

\section{Component overview}

\textbf{Contingent on Decision 1 (\S2.5).} Everything in this section
assumes \texttt{quarryd} takes \texttt{punt\_lux} as a direct, mandatory
dependency, mirroring \texttt{voxd}'s own precedent. That decision is not
yet ratified (\S2.5, Decision 1) --- if it is ruled the other way, this
component shape does not apply and needs re-deriving against whatever
constraint motivated the different ruling.

The applet is a new \texttt{QuarrySearchSubsystem}, living in
\texttt{src/quarry/daemon/}, composed into \texttt{quarryd} at startup and
run as a fire-and-forget task for the daemon's lifetime --- the same shape
as \texttt{MusicPlayerSubsystem} in \texttt{voxd} (\texttt{vox/src/}\allowbreak
\texttt{punt\_vox/voxd/music\_player/composition.py}), imported directly
and unconditionally in \texttt{voxd/daemon.py:38}, not behind a feature
flag or optional import. It owns two self-healing legs under one
\texttt{asyncio.TaskGroup}, restarting both together on any fatal fault
with backoff, exactly as \texttt{MusicPlayerSubsystem.run()} does:

\begin{description}
  \item[The push leg] Owns a \texttt{LuxClient} built for
    \texttt{quarryd}'s own app identity (\texttt{LuxClient.for\_identity},
    per \texttt{lux/src/punt\_lux/client/facade.py:76-78}) and pushes
    scene content --- the search results table, the detail view, and the
    status view --- to one named scene/frame. Modeled on
    \texttt{LuxScenePublisher}.
  \item[The receive leg] Subscribes to the events this applet's UI can
    raise (a query submission, a result selection, a tab switch between
    Search and Status) and, on each, does the corresponding work and
    re-pushes. Modeled on \texttt{LuxSubscription}/\texttt{VoxPanelLeg}'s
    \texttt{on\_event} dispatch (\texttt{vox/src/punt\_vox/panel/leg.py:}\allowbreak
    \texttt{128-130}).
\end{description}

Both legs share one \texttt{on\_connect} hook that registers the menu
entry and re-pushes whatever the subsystem last held, firing on every
fresh Hub handshake --- this is what satisfies \fr{3} (re-registration
after a \texttt{luxd} restart) for free, since it is the same mechanism
\texttt{MusicPlayerSubsystem} already uses for the identical requirement
in vox.

\section{Resolved: the list/detail swap (Open Question 1)}

\textbf{Resolved: no new board plumbing is needed.} Reading
\texttt{lux/src/punt\_lux/applets/board\_load.py} and
\texttt{board\_glass.py} in full settles this. Three facts, together:

\begin{enumerate}
  \item \texttt{BeadsBoard.request()}, called from
    \texttt{BoardLoad.board()} (\texttt{board\_load.py:73-75}), already
    returns \texttt{RenderTableRequest \textbar{} RenderRequest} for
    \emph{the same} \texttt{scene\_id}/\texttt{frame\_id} in the existing
    beads applet --- a table for issues, a plain scene for a failure or
    empty message. The Hub already composes heterogeneous request shapes
    pushed to one frame; this is not new behavior this applet would be
    the first to need.
  \item \texttt{RenderRequest.elements} (\texttt{lux/src/punt\_lux/}\allowbreak
    \texttt{operations/models/render.py:58}) is
    \texttt{list[dict[str, object]]} --- an open, self-validating element
    tree, not a short-message-only type. A full page's content, scrolled
    and rendered as text elements, is exactly the kind of thing this type
    already exists to carry.
  \item \texttt{BoardGlass.shows()} serializes pushes under one lock and
    always re-reads the held state \emph{inside} that lock before
    writing, so two pushes racing (a stale list-refresh landing after a
    fresh detail-view push, or vice versa) cannot land out of order --- the
    newer one always wins, by construction, not by a place counter that
    could itself race.
\end{enumerate}

The design, therefore: one scene id (e.g.\ \texttt{quarry-search}) serves
both the results list and the detail view. Selecting a result pushes a
fresh \texttt{RenderRequest} of the full page's content to that same
scene id, replacing the table in place; returning to the list (\fr{11})
pushes the held \texttt{RenderTableRequest} back. \texttt{BoardSlot} and
\texttt{BoardGlass} are reusable directly --- neither is
session-applet-specific, and neither's construction depends on the
\texttt{AppletProgram}/\texttt{SessionClaim} composition this applet does
not use. \texttt{AppletBoard} is \emph{not} reusable as-is: it is
transitively bound to the beads domain through its concrete
\texttt{@final BoardLoad} dependency. See ``Decisions required,'' below
--- this is a real, currently-unscoped, cross-repo piece of work, not a
free reuse.

\section{Search and selection events}

Unlike \texttt{lux-beads}, whose click carries no payload (a click just
means ``reload,'' \texttt{beads.py} takes no arguments beyond session
identity), this applet's clicks need to carry data: a query string, and a
selected result's identity (document, page, collection, chunk index)
sufficient to call \texttt{/v1/show}. \texttt{VoxPanelService}'s
\texttt{apply\_event(topic, payload)} (\texttt{vox/src/punt\_vox/panel/}\allowbreak
\texttt{service.py:146-174}) is the precedent: a payload-carrying event on
a named topic, dispatched to the field it changes, each topic named by an
enum (\texttt{PanelTopic}) rather than a bare string, so a typo in a topic
name fails at import time, not at runtime.

Two topics, minimum: \texttt{QUERY} (payload: the query string) and
\texttt{SELECT\_RESULT} (payload: enough of a result's identity to
re-fetch it). Both ride the same \texttt{listener().subscribe(*topics)}
call the push/receive legs already share, per
\texttt{VoxPanelLeg.\_listen\_once} (\texttt{leg.py:96-105}).

\section{Reaching quarry's own search}

\texttt{QuarrySearchSubsystem} runs \emph{inside} \texttt{quarryd}'s own
process. It must call the same service layer the \texttt{/v1/search} and
\texttt{/v1/show} route handlers call --- \texttt{SearchService} in
\texttt{src/quarry/daemon/routes/search.py} and its counterpart in
\texttt{src/quarry/daemon/routes/show.py}
(\texttt{src/quarry/daemon/route\_table.py} only wires route specs
to these handler methods; the handler bodies, and the service calls
worth reusing, live in \texttt{src/quarry/daemon/routes/}) --- in-process,
never looping back through its own HTTP API over loopback TLS to reach
code already running in the same process. This is a decision, not an
oversight worth
flagging as an open question: the HTTP layer's job is authenticating and
shaping requests for external clients; a subsystem composed into the
daemon itself is not an external client and gains nothing from paying
that overhead.

\section{Decisions required before implementation}

\begin{enumerate}
  \item \textbf{Packaging: \texttt{punt-lux} as a direct, mandatory
    dependency.} \texttt{voxd} imports \texttt{MusicPlayerSubsystem}
    directly and unconditionally (\texttt{daemon.py:38}), and
    \texttt{punt-lux>=0.28.0} is a plain dependency in
    \texttt{vox/pyproject.toml}, not an optional extra. The precedent is
    a hard dependency, not a feature-flagged one --- \texttt{quarryd}
    would take the same shape: \texttt{punt\_lux} always importable,
    the subsystem always composed, tolerating \texttt{luxd} being absent
    or down \emph{at runtime} (both \texttt{MusicPlayerSubsystem} and
    \texttt{VoxPanelLeg} retry forever rather than crash their host on a
    connection failure) but not absent from the dependency graph at
    install time. This needs a ratified decision before implementation,
    since it changes quarry's install footprint for every user, not only
    those who run lux.
  \item \textbf{\texttt{AppletBoard} is not actually reusable as-is; this
    needs a lux-side change.} Peer review found \S2.2's ``reuse
    \texttt{AppletBoard}/\texttt{BoardSlot}/\texttt{BoardGlass} directly''
    is imprecise. \texttt{AppletBoard.\_\_new\_\_(cls, load: BoardLoad,
    slot: BoardSlot)} (\texttt{applet\_board.py:43}) requires a concrete
    \texttt{BoardLoad} instance, and \texttt{BoardLoad} is itself
    \texttt{@final} (\texttt{board\_load.py:40}), constructed from
    concrete, beads-domain \texttt{BeadsBoard}/\texttt{BeadsSource}
    objects (\texttt{punt\_lux.apps.beads\_board}/\texttt{beads\_source}).
    \texttt{AppletBoard} is transitively bound to the beads domain through
    \texttt{BoardLoad} today. \texttt{BoardSlot} and \texttt{BoardGlass}
    genuinely have no such binding and are reusable as claimed; only
    \texttt{AppletBoard} is not. Reusing it for quarry requires a
    lux-side change --- genericizing \texttt{BoardLoad} (a Protocol, or a
    type parameter) so a \texttt{QuarrySearchBoardLoad} can satisfy
    \texttt{AppletBoard}'s constructor --- which is cross-repo work this
    document cannot unilaterally scope. Confirm with lux's own
    maintainers before assuming this reuse path is free.
  \item \textbf{App-identity client factory.} \texttt{voxd} injects its
    \texttt{LuxClient} construction behind \texttt{VoxLuxClients}
    (a factory Protocol, swapped for a fake in tests) rather than calling
    \texttt{LuxClient.for\_identity} inline. Confirm quarry's daemon
    testing conventions (\texttt{tests/hermetic\_env.py}'s embedding-fake
    pattern) extend cleanly to a second client boundary, or whether this
    needs its own hermeticity mechanism.
  \item \textbf{Status-view data source.} Per the corrected \fr{16}
    (Chapter 1), per-collection last-synced data does not exist in
    quarry's HTTP API today. Deciding whether to add it (and from what
    source --- \texttt{doctor.py}'s \texttt{compute\_sync\_plan}, or a new
    field) is a real scoping decision for whoever picks up this design,
    not resolved here.
  \item \textbf{Status-view scope beyond reachability, carried forward
    from Chapter 1.} Chapter 1's open questions ask, and this chapter
    does not resolve, exactly what \fr{13} should show beyond daemon
    reachability --- FD headroom and FTS index health are already in
    \texttt{quarry doctor}'s output and are candidates, but including
    both may be more detail than a glance-and-close status view
    warrants. Genuinely undecided; needs a ruling before
    \fr{13} can be implemented as more than a bare up/down indicator.
  \item \textbf{Success-metric instrumentation.} \S1.6's success metric
    (search volume through this applet, among users who also have
    \texttt{quarry-menubar} available) has no instrumentation point
    designed anywhere in this chapter. As currently specified, nothing
    in \texttt{QuarrySearchSubsystem} records enough to answer the
    metric's question after \pone{} ships. Before implementation, either
    design a minimal usage-logging mechanism (even a simple
    query-count-with-timestamp is likely sufficient) or explicitly
    accept the metric will need to be assessed some other way (e.g.\ a
    direct ask to users), and say so.
\end{enumerate}

\section{Known unknowns}

Board sizing and layout for a scrollable full-page detail view have not
been measured against lux's ImGui element set for a large document (a
multi-thousand-line source file, for instance) --- whether long-content
rendering performs acceptably, or needs its own pagination/virtualization
the way \fr{8} already anticipates for the results list, is untested and
should be spiked before \pone{} is considered feature-complete, not
assumed from the type system alone.

\end{document}
